% Part: computability
% Chapter: recursive-functions
% Section: pr-functions

\documentclass[../../../include/open-logic-section]{subfiles}

\begin{document}

\olfileid{cmp}{rec}{prf}
\olsection{الدوال العودية البدائية}

فلنضبط الآن الوجهين اللذين نستخرج بهما دوال جديدة من دوال معلومة:
العودية البدائية والتركيب.

\begin{defn}
  \ollabel{defn:primitive-recursion} لتكن $f$ دالة ذات $k$ حجج
  ($k\ge 1$)، ولتكن $g$ ذات $k+2$ حجة. فالمعرّفة
  \emph{بالعودية البدائية من $f$ و$g$} هي الدالة~$h$ ذات $k+1$ حجة،
  المعينة بالمعادلتين
  \begin{align*}
    h(x_0,\dots,x_{k-1},0) & =  f(x_0,\dots,x_{k-1}) \\
    h(x_0,\dots,x_{k-1},y+1) & = g(x_0,\dots,x_{k-1}, y, h(x_0,\dots,x_{k-1}, y))
  \end{align*}
\end{defn}

\begin{defn}
  \ollabel{defn:composition} لتكن $f$ ذات $k$ حجج، ولنأخذ
  $g_0$، …،~$g_{k-1}$، وهي $k$ دوال لكل منها $n$ حجة. فالدالة الحاصلة \emph{بتركيب $f$ مع $g_0$، …،~$g_{k-1}$} هي الدالة
  ~ $h$ ذات $n$ حجة التي تعرفها
  \[
  h(x_0, \dots, x_{n-1}) =
  f(g_0(x_0, \dots, x_{n-1}), \dots, g_{k-1}(x_0, \dots, x_{n-1})).
  \]
\end{defn}

ونجعل في عداد الدوال العودية البدائية، مع $\Succ$ ودوال الإسقاط
\[
\Proj{n}{i}(x_0,\dots,x_{n-1}) = x_i,
\]
لكل عدد طبيعي $n$ ولكل $i<n$، الدالة~$\Zero(x)=0$ أيضًا.

\begin{defn}
  الدوال العودية البدائية دوال من $\Nat^n$ إلى
  $\Nat$، وتُضبط مجموعتها بالتعريف الاستقرائي الآتي:
  \begin{enumerate}
  \item $\Zero$ عودية بدائية.
  \item $\Succ$ عودية بدائية.
  \item كل دالة إسقاط $\Proj{n}{i}$ عودية بدائية.
  \item إذا كانت $f$ دالة عودية بدائية ذات $k$ حجج وكانت $g_0$،
    …،~$g_{k-1}$ دوال عودية بدائية ذات $n$ حجة، فإن تركيب $f$ مع
    $g_0$، …،~$g_{k-1}$ عودي بدائي.
  \item إذا كانت $f$ دالة عودية بدائية ذات $k$ حجج وكانت $g$ دالة
    عودية بدائية ذات $k+2$ حجة، فإن الدالة المعرفة بالعودية البدائية
    من $f$ و$g$ عودية بدائية.
\end{enumerate}
\end{defn}

\begin{explain}
وحاصل التعريف أن مجموعة الدوال العودية البدائية أصغر مجموعة تضم
$\Zero$ و$\Succ$ ودوال الإسقاط~$\Proj{n}{j}$، وتكون منغلقة تحت
التركيب والعودية البدائية.

ولنا أن نصف هذه المجموعة على وجه آخر، بأن نبنيها
«مرحلة» بعد مرحلة. فلتكن $S_0$ مجموعة الدوال التي نبدأ بها، وهي $\Zero$ و$\Succ$ ودوال
الإسقاط. فهذه دوال المرحلة~$0$. فإذا فرغنا من
تعريف $S_i$، جعلنا $S_{i+1}$ مجموعة ما يحصل بإجراء
تركيب واحد أو عودية بدائية واحدة على دوال من~$S_i$.
فيكون
\[
S = \bigcup_{i \in \Nat} S_i
\]
مجموعة جميع الدوال العودية البدائية.
\end{explain}

ولنطبّق هذا الضبط فنثبت أن $\Add$ دالة عودية بدائية.

\begin{prop}
  دالة الجمع $\Add(x,y)=x+y$ عودية بدائية.
\end{prop}

\begin{proof}
قد سبق تعريف $\Add$ بالعودية البدائية من دالتين $f$ و~$g$،
وجاء التعريف على صورة \olref{defn:primitive-recursion}:
\begin{align*}
  \Add(x_0, 0) & = f(x_0) = x_0\\
  \Add(x_0, y+1) & = g(x_0, y, \Add(x_0, y)) =
  \Succ(\Add(x_0, y))
\end{align*}
فلا يبقى لإثبات عودية $\Add$ البدائية إلا إثباتها لـ$f$ و$g$. أما الأولى فإن
$f(x_0)=x_0=\Proj{1}{0}(x_0)$، ودوال الإسقاط تعد عودية بدائية، فتكون
$f$ عودية بدائية. أما $g$ فهي الدالة الثلاثية الحجج $g(x_0,y,z)$
المعرفة بـ
\[
g(x_0, y, z) = \Succ(z).
\]
ولا يكفي هذا وحده للحكم بأن $g$ عودية بدائية؛ فإن $g$ غير $\Succ$
من جهة عدد الحجج: فـ$\Succ$ أحادية، و$g$ لا بد أن تكون ثلاثية
الحجج. وإنما يتم تعريف $g$ على الوجه الرسمي بالتركيب الآتي:
\[
g(x_0, y, z) = \Succ(\Proj{3}{2}(x_0, y, z)).
\]
وحيث إن $\Succ$ و$\Proj{3}{2}$ عوديتان بدائيتان، فإن $g$
عودية بدائية أيضًا، لكونها تركيبًا منهما.
\end{proof}

\begin{prop}
  \ollabel{prop:mult-pr}
  دالة الضرب $\Mult(x,y)=x\cdot y$ عودية بدائية.
\end{prop}

\begin{proof}
  تمرين.
\end{proof}

\begin{prob}
  برهن \olref[cmp][rec][prf]{prop:mult-pr}: ضع تعريف $\Mult$
  بالعودية البدائية في الصورة المشترطة في
  \olref[cmp][rec][prf]{defn:primitive-recursion}، ثم بيّن أن الدالتين
  اللتين تقومان فيه مقام $f$ و~$g$ عوديتان بدائيتان.
\end{prob}

\begin{ex}
ولنعد إلى أول ما أوردناه مثالًا للتعريف بالعودية البدائية:
\begin{align*}
h(0) & =  1 \\
h(y+1) & =  2 \cdot h(y).
\end{align*}
لا يندرج هذا التعريف كما هو في الصورة المشترطة في
\olref{defn:primitive-recursion}، لأن $k=0$. ويتضمن التعريف أيضًا
الثابتين $1$ و$2$. فأما الإشكال الأول فندفعه بإدخال حجة صورية، فنعرّف
الدالة~$h'$ هكذا:
\begin{align*}
h'(x_0, 0) & =  f(x_0) = 1 \\
h'(x_0, y+1) & =  g(x_0, y, h'(x_0, y)) = 2 \cdot h'(x_0, y).
\end{align*}
يمكن تعريف الدالة $f(x_0)=1$ من $\Succ$ و$\Zero$ بالتركيب:
$f(x_0)=\Succ(\Zero(x_0))$. ويمكن تعريف الدالة $g$ بالتركيب من
$g'(z)=2\cdot z$ ومن دوال الإسقاط:
\begin{align*}
g(x_0, y, z) & = g'(\Proj{3}{2}(x_0, y, z))
\intertext{ويمكن بدورها تعريف $g'$ بالتركيب هكذا:}
g'(z) & = \Mult(g''(z), \Proj{1}{0}(z))
\intertext{حيث}
g''(z) & = \Succ(f(z)),
\end{align*}
و$f$ كما سبق: $f(z)=\Succ(\Zero(z))$. وحيث حصلت لنا~$h'$،
فلنستعمل التركيب مرة أخرى ونضع
$h(y)=h'(\Proj{1}{0}(y),\Proj{1}{0}(y))$. فقد استخرجنا $h$
من الدوال الأساسية بتركيبات وعوديات بدائية متتابعة؛ وبذلك ثبت أن
$h$ عودية بدائية.
\end{ex}

\end{document}
